% Source: https://artofproblemsolving.com/community/c3h1220659_hexagons
% Stable identity: aops-topic-1220659-q02
\textbf{No.} Take an equilateral triangle of side $4$ and cut off an equilateral triangle of side $1$ at each corner. The remaining hexagon has every angle equal to $120^\circ$, but its sides alternate $2$ and $1$.

\begin{asy}
usepackage("amsmath");
size(3cm);
defaultpen(fontsize(10pt));
pair A=(0,0),B=(4,0),C=(2,2*sqrt(3)); draw(A--B--C--cycle,gray+dashed);
pair P1=A+(B-A)/4,P2=B-(B-A)/4,P3=B+(C-B)/4,P4=C-(C-B)/4,P5=C+(A-C)/4,P6=A-(A-C)/4;
draw(P1--P2--P3--P4--P5--P6--cycle,blue+linewidth(1)); label("$2$",(P1+P2)/2,S); label("$1$",(P2+P3)/2,SE);
label("\mbox{}",(min(currentpicture,user=true)+max(currentpicture,user=true))/2,black); // keeps the text after the figure black when LaTeX typesets the labels inline
\end{asy}

For triangles the situation differs: equal angles force equal sides because side lengths are determined by the angles up to scale (law of sines). For polygons with more sides, the angles do not determine the shape.
